%% ma: L12-B14-0132
%% lop: 12
%% chuong: 5
%% bai: 14
%% dang: 12.14.21
%% loai: TN4
%% muc: TH
%% dap_an: "C"
%% nguon: "Ảnh giáo viên gửi 25/11/2026 (tài liệu ôn tập Phương trình mặt phẳng), Phần I Câu 23"
%% kien_thuc: "Bài 14 (phương trình mặt phẳng theo đoạn chắn), Bài 7 (trọng tâm tam giác)"
%% trang_thai: chinh-thuc
\begin{ex}
Trong không gian $Oxyz$, mặt phẳng $(P)$ cắt $3$ trục tọa độ tại ba điểm phân biệt tạo thành một tam giác có trọng tâm $G(3;2;-1)$. Viết phương trình mặt phẳng $(P)$.
\choice{$\dfrac x9+\dfrac y6+\dfrac z3=1$}{$\dfrac x9+\dfrac y6-\dfrac z3=0$}{\True $\dfrac x9+\dfrac y6-\dfrac z3=1$}{$\dfrac x9+\dfrac y6+\dfrac z3=0$}
\loigiai{Gọi $A(a;0;0)$, $B(0;b;0)$, $C(0;0;c)$. Vì $G$ là trọng tâm nên $\dfrac a3=3$, $\dfrac b3=2$, $\dfrac c3=-1$, suy ra $a=9$, $b=6$, $c=-3$. Phương trình $(P)$: $\dfrac x9+\dfrac y6+\dfrac z{-3}=1$.}
\end{ex}
